% ======================================================================
%  particulars/software.tex — CONTENT for the Software resume
% ----------------------------------------------------------------------
%  This file only defines content macros; it does no formatting. main.tex
%  renders these macros into the final layout.
%
%  To make a new variant (e.g. for a different industry or company):
%    1. Copy this file to a new name, e.g. particulars/swe.tex
%    2. Edit the content below.
%    3. Point main.tex's \input line at your new file.
%
%  The macros expected by main.tex are:
%    \firstname, \lastname, \email, \linkedin, \phone   (header)
%    \summaryContent   (optional - shown only if uncommented in main.tex)
%    \skillsContent    (Skills section)
%    \educationEntries (Education section)
%    \experienceEntries(Experience section)
%    \projectsEntries  (Projects section)
% ======================================================================

% ----------------------------------------------------------------------
%  Personal details
% ----------------------------------------------------------------------
\def\firstname{Chiedozie}
\def\lastname{Ihebuzor}
\def\email{your.email@domain.com}
\def\linkedin{www.linkedin.com/in/chiedozie-ihebuzor}
\def\phone{+44 1234 56789}

% ----------------------------------------------------------------------
%  Summary (optional)
% ----------------------------------------------------------------------
\def\summaryContent{
\section{Summary}
Passionate, hard-working engineer with a love for learning. Eager to take responsibility and build solutions that make a real impact within the team and company. Uses data to guide priorities and decisions, and leverages strong communication skills to learn from other teams and share insights across the organisation.}

% ----------------------------------------------------------------------
%  Skills
% ----------------------------------------------------------------------
\def\skillsContent{
\begin{flushleft}
    C++ \textbar\ 
    Python \textbar\ 
    Git \textbar\ 
    Linux \textbar\ 
    Bash \textbar\ 
    G-Test \textbar\
    SCons \textbar\ 
    CMake \textbar\ 
    Networks \textbar\
    Hardware Design \textbar\ 
    FPGAs \textbar\
    Verilog \textbar\
\end{flushleft}
}

% ----------------------------------------------------------------------
%  Education
% ----------------------------------------------------------------------
%  Each entry: \resumeSubheading{school}{dates}{course, grade}{} followed
%  by an optional bullet list.
\def\educationEntries{
    \resumeSubheading{Imperial College London}{October 2022 -- June 2026}
        {Electrical and Electronic Engineering, Current Grade: First Class Honours}{}
    \resumeItemListStart
        \resumeItem{Key Modules: Digital System Design, Finance, Computer Architecture, ML, Deep Learning, VLSI (Apple-Sponsored).}
        \resumeItem{Societies: Robotics Society, Live Music Society (Guitar Technician), Skateboarding, Cheerleading, Gospel Choir.}
    \resumeItemListEnd

    \resumeSubheading{King Edward IV College}{September 2020 -- May 2022}
        {A-Levels, Physics (A*), Computer Science (A*), Maths (A*), AS Further Maths (A)}{}
}

% ----------------------------------------------------------------------
%  Experience
% ----------------------------------------------------------------------
%  Each entry: \resumeSubheading{role}{dates}{company}{location} followed
%  by a bullet list of highlights.
\def\experienceEntries{
    \resumeSubheading{Quantitative FPGA Technology Placement}{Mar 2025 -- Sept 2025}{Qube Research \& Technologies}{London}
    \resumeItemListStart
        \resumeItem{Designed a high-throughput, low-latency ML inference engine on FPGA to accelerate and up-scale existing software trading decision model.}
        \resumeItem{Collaborated across quantitative development, research and infrastructure teams to migrate and optimise the trading model for FPGA deployment.}
        \resumeItem{Optimised C++ API to write model features to FPGA over low-latency network card with minimum latency.}
        \resumeItem{Parallelised team’s hardware unit simulation CI with Python script, managing licences across multiple hosts and processes, reducing max runtime by 8x.}
    \resumeItemListEnd
    \resumeSubheading{Quantitative Infrastructure Technology Intern}{Jun 2024 -- Sept 2024}{Qube Research \& Technologies}{London}
    \resumeItemListStart
        \resumeItem{Developed terminal-based C++ monitoring tool for L3 cache performance of trading and market access systems.}
        \resumeItem{Simplified company-wide cache policy allocation process with a terminal UI, reducing set-up errors}
        \resumeItem{Automated FPGA diagnostics logging by parsing hardware register maps and generating C++ access code via jinja templating.}
    \resumeItemListEnd
    \resumeSubheading{Research Intern}{Aug 2023 -- Sept 2023}{Imperial College London - Department of Bioengineering}{London}
    \resumeItemListStart
        \resumeItem{Designed and developed a modular arcade platform to support game-based science outreach programs.}
        \resumeItem{Produced detailed documentation for ease of use by other researchers.}
    \resumeItemListEnd
    \resumeSubheading{Private Tutor}{Sept 2022 -- Present}{Self-Employed}{Remote}
    \resumeItemListStart
        \resumeItem{Delivered tailored tutoring in A-Level Mathematics and Physics to students to students aiming for top grades.}
    \resumeItemListEnd
}

% ----------------------------------------------------------------------
%  Projects
% ----------------------------------------------------------------------
%  Each entry: \resumeProjectHeading{\textbf{title}}{} followed by a bullet
%  list, then a small negative \vspace to tighten the gap before the next
%  project.
\def\projectsEntries{
    \resumeProjectHeading{\textbf{FPGA CORDIC Accelerator}}{}
    \resumeItemListStart
        \resumeItem{Implemented an unrolled, pipelined CORDIC module on an Altera FPGA to accelerate trigonometric evaluations.}
        \resumeItem{Benchmarked the module against standard C libraries to show significant speedups in calculations.}
    \resumeItemListEnd
    \vspace{-13pt}
    \resumeProjectHeading{\textbf{Skateboarding Data Tracker}}{}
    \resumeItemListStart
        \resumeItem{Wrote real-time trick detection algorithm using accelerometer, gyroscope, and ultrasonic sensor data over I\textsuperscript{2}C.}
        \resumeItem{Streamed motion analytics to Firebase and displayed data on a custom front-end dashboard.}
    \resumeItemListEnd
    \vspace{-10pt}
    \resumeProjectHeading{\textbf{Rubik’s Cube Solver}}{}
    \resumeItemListStart
        \resumeItem{Implemented a solving algorithm in Python, featuring an interactive 3D visualization with real-time mouse control.}
        \resumeItem{Used OpenCV to import a user’s physical cube state for automated step-by-step solving.}
    \resumeItemListEnd
    \vspace{-10pt}
}
